% Part: computability
% Chapter: computability-theory
% Section: s-m-n

\documentclass[../../../include/open-logic-section]{subfiles}

\begin{document}

\olfileid{cmp}{thy}{smn}
\olsection{$s$-$m$-$n$ 定理}

\begin{explain}
下一个定理被称为“$s$-$m$-$n$定理”，稍后就会明白其得名原因。
理解这个定理究竟在说什么才是难点所在；一旦理解了它的内容，就会发现它其实相当显然。
\end{explain}

\begin{thm}
\ollabel{thm:s-m-n}
  对于任意一对自然数 $n$ 和~$m$，都存在一个原始递归函数~$s^m_n$，
  使得对任意 $e$, $a_0$, \dots, $a_{m-1}$, $y_0$, \dots, $y_{n-1}$，都有
  \[
  \cfind{s^m_n(e, a_0, \dots, a_{m-1})}[n](y_0, \dots, y_{n-1}) \simeq
  \cfind{e}[m+n](a_0, \dots, a_{m-1}, y_0, \dots, y_{n-1}).
\]
\end{thm}

\begin{explain}
将 $s^m_n$ 视为对\emph{程序}的操作会很有助于理解。
也就是说，$s^m_n$ 以一个 $(m+n)$ 元函数的程序~$e$ 以及固定的输入 $a_0$, \dots, $a_{m-1}$ 为自变元，并返回一个程序
$s^m_n(x, a_0, \dots, a_{m-1})$，该程序对应于由剩余自变元构成的 $n$ 元函数。\iftag{TMs}{若将 $e$ 视为一台图灵机的描述，那么 $s^m_n(e, a_0, \dots, a_{m-1})$ 便是这样一台图灵机：当输入为 $y_0$, \dots,~$y_{n-1}$ 时，它将 $a_0$, \dots,~$a_{m-1}$ 前置到输入串上，并运行~$e$。由此，每个 $s^m_n$ 不过是一个寻找相应图灵机编码的原始递归函数。}{}
\end{explain}

\end{document}